\scanpage{21}
\mspage{20}

$B$ sur $S$ et un hom de groupes

\[
B\xrightarrow{i}\underline{\mathrm{Pic}}_{X/S},
\]

ce qui définit d'ailleurs un torseur $Q$ sous $A=\underline{\mathrm{Pic}}^0_{B/S}$ et un hom canonique

\[
X\xrightarrow{\psi} Q
\]

[C'est par un \unsure{sorite|suite} sur la correspondance entre des $i$ et des $\psi$ que devrait commencer \unsure{mon|un|une} \deleted{\unsure{le}} petit topo de Albanese\ldots.].

Tout ce qui a été dit \unsure{vaudra|vaut} alors, on trouve

\[
X\times\underline{\mathrm{Pic}}_{X/S}\mathbin{\unsure{-|\to}}\underline{\mathrm{Pic}}_{B/S}
\]

définissant

\[
X\times\underline{\mathrm{Pic}}^\delta_{X/S}\longrightarrow\underline{\mathrm{Pic}}^{\varphi(\delta)}_{B/S}
\]

\margin{\uline{NB} : \unsure{sous-entendus|sous-entendre} dans la notation $\varphi(\delta)$}

qui donne, en fait, si on a un $B$\deleted{\unsure{-teur|-torseur}} $P$ et un hom $P\longrightarrow\underline{\mathrm{Pic}}^\delta_{X/S}$ \inserted{compatible avec $i$}

\[
X\longrightarrow\underline{\mathrm{Pic}}^{\varphi(\delta)}_{B/S}\overset{A}{\times}\bigl(\deleted{\underline{\mathrm{Pic}}_{X/S}^{\delta}}P\overset{B}{\times}(A,\widetilde{\varphi(\delta)})\bigr)
\]

\deleted{\unsure{et}}\inserted{se factorisant par}

\[
Q^{-\chi(\delta)}\longrightarrow\underline{\mathrm{Pic}}^{\varphi(\delta)}_{B/S}\overset{A}{\times}\bigl(\deleted{\underline{\mathrm{Pic}}^\delta_{X/S}}P\overset{B}{\times}(A,\widetilde{\varphi(\delta)})\bigr)
\]

i.e. on a une section de

\[
Q^{\chi(\delta)}\overset{A}{\times}\underline{\mathrm{Pic}}^{\varphi(\delta)}_{B/S}\overset{A}{\times}\bigl(\deleted{\underline{\mathrm{Pic}}^\delta_{X/S}}P\overset{B}{\times}(A,\widetilde{\varphi(\delta)})
\]

\deleted{b) On peut aller plus loin, dans la \deleted{\unsure{réduction}} simplification} \deleted{et généralisation des énoncés, lorsqu'on se borne à} \deleted{des $X$ munis de sections \unsure{marquées} $g$, \deleted{Alors} du} \deleted{fait qu'on a $X\longrightarrow A$. \unsure{Ceci noté}, $X$ n'interv}
